\section*{Hoofdstuk \ref{ch:b2:liquid-vapour-diagrams} --- Binaire vloeistof-dampdiagrammen en destillatie}

\begin{solution}{exo:b2:liquid-vapour-diagrams:1}
Benzeen \qty{80.1}{\degreeCelsius}, tolueen \qty{110.6}{\degreeCelsius}. De vloeistof
met samenstelling 0.2 begint te koken bij ongeveer \qty{102}{\degreeCelsius}; haar eerste
damp heeft samenstelling 0.38.
\end{solution}

\begin{solution}{exo:b2:liquid-vapour-diagrams:2}
$p = 0.5 \times 1.010 + 0.5 \times 0.388 = \qty{0.699}{bar}$;
$y = 0.505/0.699 = 0.722$.
\end{solution}

\begin{solution}{exo:b2:liquid-vapour-diagrams:3}
$n_V/n = (0.30 - 0.20)/(0.45 - 0.20) = 0.40$: \qty{4.0}{mol} damp en
\qty{6.0}{mol} vloeistof. Controle: $6.0 \times 0.20 + 4.0 \times 0.45 = 3.0 = 10
\times 0.30$.
\end{solution}

\begin{solution}{exo:b2:liquid-vapour-diagrams:4}
(a) $v = 2 + 2 - 1 = 3$: $T$, $p$ en $x$ zijn vrij. (b) $v = 2 + 2 - 2 = 2$; bij
vaste druk 1. (c) De \omterm{def:b2:liquid-vapour-diagrams:azeotrope}{azeotroop} voegt de betrekking $x = y$ toe: $v = 1$; bij vaste
druk 0 --- de \omterm{def:b2:liquid-vapour-diagrams:azeotrope}{azeotroop} kookt bij een vaste temperatuur, zoals een zuivere vloeistof,
al verandert zijn samenstelling met de druk.
\end{solution}

\begin{solution}{exo:b2:liquid-vapour-diagrams:5}
Bij het \omterm{def:b2:liquid-vapour-diagrams:bubble-dew}{dauwpunt} is $yp/p_1^* + (1 - y)p/p_2^* = 1$. Bij \qty{100}{\degreeCelsius}:
$0.3 \times 1.013/1.80 + 0.7 \times 1.013/0.742 = 1.124 > 1$ (te koud); bij
\qty{105}{\degreeCelsius}: $0.148 + 0.824 = 0.971 < 1$. Interpoleren geeft $T \approx
100 + 5 \times 0.124/0.153 = \qty{104}{\degreeCelsius}$. De eerste druppel heeft $x =
yp/p_1^* \approx 0.3 \times 1.013/2.0 = 0.15$.
\end{solution}

\begin{solution}{exo:b2:liquid-vapour-diagrams:6}
De eerste druppels hebben de samenstelling van de damp boven de kokende vloeistof,
$y = 0.71$, bij \qty{92.1}{\degreeCelsius}. Omdat benzeen bij voorkeur vertrekt, wordt de
vloeistof er armer aan, stijgt haar kookpunt en wordt ook het destillaat
armer; de laatste druppels zijn bijna zuiver tolueen. Eenvoudige destillatie is één
enkele evenwichtstrap, herhaald op een vloeistof die voortdurend verandert: elke
fractie is een mengsel.
\end{solution}

\begin{solution}{exo:b2:liquid-vapour-diagrams:7}
Voeding 0.05, aan de waterkant van de \omterm{def:b2:liquid-vapour-diagrams:azeotrope}{azeotroop}: de top levert de \omterm{def:b2:liquid-vapour-diagrams:azeotrope}{azeotroop}
(0.88 in mol, 95~\% in massa), de ketel houdt water. Voeding 0.95, aan de
ethanolkant: de top levert nog steeds de \omterm{def:b2:liquid-vapour-diagrams:azeotrope}{azeotroop}, het laagste kookpunt,
en de ketel houdt zuivere ethanol.
\end{solution}

\begin{solution}{exo:b2:liquid-vapour-diagrams:8}
Het mengsel kookt wanneer $0.10 + p^*_{\text{water}} = 1.013$, dus
$p^*_{\text{water}} = \qty{0.91}{bar}$, bij ongeveer \qty{97}{\degreeCelsius}. Massaverhouding
$0.10 \times 136/(0.91 \times 18.02) = 0.83$: \qty{0.83}{g} essence per gram
water.
\end{solution}

\begin{solution}{exo:b2:liquid-vapour-diagrams:9}
Bij kamertemperatuur twee lagen L$_2$ en L$_1$ (samenstelling 0.4 ligt in het
\omterm{def:b2:liquid-vapour-diagrams:miscibility-gap}{ontmengingsgebied}). Bij verwarmen koken de lagen samen bij de temperatuur van de
\omterm{def:b2:liquid-vapour-diagrams:miscibility-gap}{heteroazeotroop}, die constant blijft; de eerste damp heeft de heteroazeotrope
samenstelling (het punt). Omdat 0.4 aan de kant van L$_2$ van het punt ligt, wordt de laag
L$_1$ eerst opgebruikt; daarna stijgt de temperatuur, wordt de resterende L$_2$
armer aan 1 langs haar \omterm{def:b2:liquid-vapour-diagrams:bubble-dew}{kookkromme} en volgt de damp de \omterm{def:b2:liquid-vapour-diagrams:bubble-dew}{dauwkromme}, tot
het punt de \omterm{def:b2:liquid-vapour-diagrams:bubble-dew}{dauwkromme} bereikt: de laatste druppel verdampt.
\end{solution}

\begin{solution}{exo:b2:liquid-vapour-diagrams:10}
$N_{\min} = \ln(99 \times 99)/\ln 2.47 = 9.19/0.904 = 10.2$, tegenover 6.5: van
95~\% naar 99~\% zuiverheid aan beide uiteinden kost meer dan de helft
schotels extra. Elke extra factor op $x/(1 - x)$ kost hetzelfde aantal schotels, en
zuiverheid wordt gemeten aan de overblijvende onzuiverheid, die slechts meetkundig daalt.
\end{solution}

\begin{solution}{exo:b2:liquid-vapour-diagrams:11}
Het \omterm{def:b2:liquid-vapour-diagrams:binary-diagram}{isobare diagram} alleen geeft het niet: men heeft $p_1^*$ en $p_2^*$ nodig bij
\qty{80}{\degreeCelsius}, uit de vergelijkingen van Antoine (1.010 en
\qty{0.388}{bar}). Dan is $p = p_2^* + x(p_1^* - p_2^*)$ en $p = 1/(y/p_1^* + (1 -
y)/p_2^*)$. De vloeistof is stabiel bij hoge druk (samenpersen van een damp
condenseert hem), en bij lage temperatuur in het \omterm{def:b2:liquid-vapour-diagrams:binary-diagram}{isobare diagram}: de twee diagrammen
staan ten opzichte van elkaar op hun kop.
\end{solution}

\begin{solution}{exo:b2:liquid-vapour-diagrams:12}
Omdat $y$ toeneemt met $x$ en $y(1 - y) > 0$, heeft $dp/dx$ het teken van $y - x$:
de druk stijgt met $x$ waar de damp rijker is aan 1. Voor een \omterm{def:b2:chemical-potential:ideal-mixture}{ideaal
mengsel} is $y > x$ voor elke $0 < x < 1$ (\cref{prop:b2:liquid-vapour-diagrams:richer}):
$p$ is strikt monotoon, heeft geen extremum, en er is geen \omterm{def:b2:liquid-vapour-diagrams:azeotrope}{azeotroop}.
\end{solution}

\begin{solution}{pb:b2:liquid-vapour-diagrams:1}
\textbf{1.} $T = B/(A - \log_{10}1.013) - C$: benzeen $1203.835/4.01253 + 53.226 =
\qty{353.2}{K}$, \qty{80.1}{\degreeCelsius}; tolueen $1343.943/4.07266 + 53.773 =
\qty{383.8}{K}$, \qty{110.6}{\degreeCelsius}.
\textbf{2.} Bij \qty{363.15}{K}: $p_1^* = \qty{1.361}{bar}$, $p_2^* = \qty{0.542}{bar}$.
\textbf{3.} $x = (1.013 - 0.542)/(1.361 - 0.542) = 0.575$; $y = 0.575 \times
1.361/1.013 = 0.773$.
\textbf{4.} $x = (1.013 - 0.742)/(1.800 - 0.742) = 0.256$; $y = 0.256 \times
1.800/1.013 = 0.455$.
\textbf{5.} Kookpunten (0, 110.6), (0.256, 100), (0.575, 90), (1, 80.1); \omterm{def:b2:liquid-vapour-diagrams:bubble-dew}{dauwpunten}
(0, 110.6), (0.455, 100), (0.773, 90), (1, 80.1) --- de lens van de
figuur.
\textbf{6.} $\frac12(p_1^* + p_2^*) = 1.013$: bij \qty{90}{\degreeCelsius} 0.952, bij
\qty{95}{\degreeCelsius} 1.102; interpoleren geeft $90 + 5 \times 0.061/0.150 \approx
\qty{92}{\degreeCelsius}$ (precies \qty{92.1}{\degreeCelsius}).
\textbf{7.} $v = 2$, 1 bij vaste druk: elke temperatuur legt $x$ en $y$ vast,
en daarom bestaat het diagram uit twee krommen.
\textbf{8.} $x = (1.013 - 0.636)/(1.569 - 0.636) = 0.404$; $y = 0.404 \times
1.569/1.013 = 0.626$.
\textbf{9.} $n_V = 100 \times (0.500 - 0.404)/(0.626 - 0.404) = \qty{43}{mol}$,
$n_L = \qty{57}{mol}$.
\textbf{10.} Damp $43 \times 0.626 = \qty{26.9}{mol}$; vloeistof $57 \times 0.404
= \qty{23.0}{mol}$; totaal \qty{49.9}{mol} $\approx 50$.
\textbf{11.} $26.9/50 = 54~\%$ van het benzeen, in een damp met slechts 63~\%.
\textbf{12.} Bij \qty{100}{\degreeCelsius} heeft de damp in evenwicht $y =
0.455 < 0.5$: het voedingspunt ligt boven de \omterm{def:b2:liquid-vapour-diagrams:bubble-dew}{dauwkromme}, de voeding is volledig
damp en er wordt niets gescheiden.
\textbf{13.} Bij haar \omterm{def:b2:liquid-vapour-diagrams:bubble-dew}{dauwpunt}, \qty{98.8}{\degreeCelsius} (met de methode van
\cref{exo:b2:liquid-vapour-diagrams:5}).
\textbf{14.} $\alpha = p_1^*/p_2^*$; bij \qty{80.1}{\degreeCelsius} is dat $1.013/0.390 =
2.60$.
\textbf{15.} Bij \qty{110.6}{\degreeCelsius}: $2.377/1.013 = 2.35$.
\textbf{16.} $\sqrt{2.60 \times 2.35} = 2.47$.
\textbf{17.} $y = xp_1^*/p$ en $1 - y = (1 - x)p_2^*/p$; deel op elkaar.
\textbf{18.} Er wordt niets afgetapt: tussen twee schotels voeren de opstijgende damp en
de neerdalende vloeistof dezelfde hoeveelheid materie en van elk bestanddeel,
$V y_k = L x_{k+1}$ met $V = L$, dus $x_{k+1} = y_k$.
\textbf{19.} Met $r = x/(1 - x)$: $r(y_k) = \alpha\,r(x_k)$ en $x_{k+1} = y_k$, dus
$r(x_{k+1}) = \alpha\,r(x_k)$ en $r(x_D) = \alpha^N r(x_B)$; $N = \ln[r(x_D)/
r(x_B)]/\ln\alpha$.
\textbf{20.} $N_{\min} = \ln(19 \times 19)/\ln 2.47 = 5.889/0.904 = 6.5$.
\textbf{21.} Een kolom moet product leveren, dus werkt ze bij een eindige \omterm{def:b2:liquid-vapour-diagrams:fractional-distillation}{refluxverhouding},
wat meer schotels vraagt; en een echte schotel bereikt het evenwicht niet (zijn
rendement is kleiner dan één).
\textbf{22.} Zeven treden, de laatste eindigend bij 0.034, onder 0.05: 6.5 schotels in een
continue telling, zeven hele (de verdamper telt als één).
\textbf{23.} $x = (95/46.07)/(95/46.07 + 5/18.02) = 0.881$.
\textbf{24.} Bovenaan op zijn best de \omterm{def:b2:liquid-vapour-diagrams:azeotrope}{azeotroop} (0.88 in mol); onderaan
water. Vijftig schotels steken de \omterm{def:b2:liquid-vapour-diagrams:azeotrope}{azeotroop} niet over.
\textbf{25.} Bij totale reflux $\boldsymbol{N_{\min} \approx 6.5}$ \omterm{def:b2:liquid-vapour-diagrams:fractional-distillation}{theoretische
schotels}.
\end{solution}
